\documentclass[11pt]{article}
\usepackage{ap-lang}
\title{Part 2: Claims and Evidence: Question Set}
\author{Paul Wang}
\date{November 2026}
\docpart{Part 2 \textperiodcentered\ Question Set}

\begin{document}
\maketitle

\begin{abstract}
Three exam-style passages with multiple-choice questions, followed by a full argument essay (Question 3). Passage 1 is an essay arguing that the paths people wear across lawns are evidence planners should heed; Passage 2 is an op-ed against minimum parking requirements. Both are built on claims, evidence, and concession, so most questions ask what a piece of evidence does or how a claim is qualified. Passage 3 is a writing set on a student's argument for a personal finance requirement. All passages, writers, places, and figures are fictional and were written for this practice set. Allow about 35 minutes for the multiple choice and 40 minutes for the essay.
\end{abstract}

\scoringline{MCQ: \blank\ / \qtotal \qquad Essay: \blank\ / 6 \qquad Target: 19+ MCQ and 4+ essay \qquad Pace: 80 s per MCQ}

\section{Passage 1: Desire paths (reading)}

\begin{passage}{Passage 1}{The following passage is from an essay published in 2024 in a magazine about cities and design.}
\pn{1} Every campus has one: a strip of bare dirt cutting diagonally across a lawn that the architects intended to be crossed only along its edges. Groundskeepers call it a ``desire path,'' and they usually treat it as vandalism. They put up a little fence. They plant a hedge. They post a sign that says, with touching optimism, PLEASE USE SIDEWALK.

\pn{2} I want to suggest that the dirt is not the problem. The dirt is data.

\pn{3} A desire path is the most honest survey a planner will ever receive. No one fills it out to please the researcher; no one exaggerates. It records, one footstep at a time, where people actually need to go, as opposed to where someone at a drafting table imagined they would. When thousands of people independently choose the same line, they are not being careless. They are voting.

\pn{4} Some designers have learned to count those votes. When Halvorsen College built its new science quad in 2009, its planners did something unusual: they planted grass everywhere and laid no paths at all for the first year. By spring, students had worn a web of trails across the lawn. The college paved exactly those trails. Fifteen years later, the quad has needed almost no fencing and almost no reseeding, and its maintenance budget is a third of that of the older quad next door.

\pn{5} Of course, not every desire path deserves pavement. Some cross wetlands that took decades to recover from the first few hundred shortcuts; some lead people across traffic where they should not go. A path that is popular is not thereby safe. But the right response to a dangerous desire path is to ask what need it reveals---a bus stop too far from the crosswalk, a gate locked at the wrong hour---and to meet that need somewhere safer. The fence answers the symptom. The question answers the cause.

\pn{6} The lesson reaches well beyond lawns. Institutions of every kind design for the people they imagine and then scold the people who actually show up. A tax form that half its users fill out wrong is not evidence of a careless public; a medication schedule that patients routinely abandon is not evidence of irresponsible patients. Before we build another fence, we might look at where the grass has worn thin, and ask what it has been trying to tell us.
\end{passage}

\begin{q}{4}{2.4 Counterargument}{Paragraph 1 presents, with gentle irony, the conventional response the writer will challenge.}
In paragraph 1, the description of how groundskeepers respond to desire paths primarily serves to
\ch{praise groundskeepers for protecting the lawn}{explain how desire paths are formed}{provide evidence that desire paths damage campuses}{present, with gentle irony, the conventional view the writer goes on to challenge}
\end{q}

\begin{q}{3}{2.1 Claims}{``The dirt is data'' compactly states the central claim.}
Paragraph 2 functions primarily as
\ch{an anecdote that illustrates a later claim}{a transition to a new subject}{a compact statement of the writer's central claim}{a concession to the groundskeepers' view}
\end{q}

\begin{q}{4}{4.3 Comparisons}{Calling the path a survey and a vote recasts walkers' behavior as meaningful evidence.}
In paragraph 3, the writer's comparisons of a desire path to a ``survey'' and to ``voting'' primarily serve to
\ch{criticize students for ignoring campus rules}{show that desire paths are a recent phenomenon}{suggest that planners should conduct more surveys}{recast the walkers' behavior as meaningful collective evidence rather than carelessness}
\end{q}

\begin{q}{1}{2.2 Evidence}{The Halvorsen example shows the writer's approach producing practical benefits.}
The example of Halvorsen College (paragraph 4) functions mainly as
\ch{evidence that following desire paths can produce practical benefits}{a counterexample that complicates the writer's claim}{expert testimony from the college's planners}{a personal anecdote about the writer's education}
\end{q}

\begin{q}{1}{2.2 Evidence}{The budget comparison is the detail that measures resources saved.}
Which detail from paragraph 4 most directly supports the idea that designing around desire paths saves resources?
\ch{``its maintenance budget is a third of that of the older quad next door''}{``students had worn a web of trails across the lawn''}{``When Halvorsen College built its new science quad in 2009''}{``its planners did something unusual''}
\end{q}

\begin{q}{3}{2.4 Qualification}{The writer limits the claim (not every path) and redirects to the need behind it.}
Paragraph 5 primarily serves to
\ch{argue that all fences should be removed}{abandon the argument that planners should follow desire paths}{qualify the writer's claim by acknowledging its limits, then redirect attention to the need a path reveals}{introduce statistics about wetland recovery}
\end{q}

\begin{q}{2}{4.2 Syntax}{Parallel sentences contrast the two responses to show why asking is better.}
The sentences ``The fence answers the symptom. The question answers the cause.'' (paragraph 5) use parallel structure mainly to
\ch{summarize the Halvorsen College example}{contrast two responses and show why the writer's is more effective}{introduce a new counterargument}{suggest that fences and questions are equally useful}
\end{q}

\begin{q}{4}{3.3 Structure}{The final paragraph extends the principle beyond lawns, broadening its significance.}
The examples of the tax form and the medication schedule in paragraph 6 primarily serve to
\ch{shift the essay's focus from design to medicine}{provide statistical support for the Halvorsen example}{concede that some users are careless}{extend the writer's principle to other institutions, broadening its significance}
\end{q}

\begin{q}{2}{2.1 Claims}{The argument rests on the assumption that widespread voluntary behavior reveals real needs.}
The writer's argument relies most heavily on which assumption?
\ch{Paved paths are always cheaper to maintain than grass.}{When many people independently choose the same behavior, it reveals a genuine need that design should accommodate.}{Groundskeepers are generally careless at their jobs.}{All shortcuts are safe if enough people take them.}
\end{q}

\section{Passage 2: Parking minimums (reading)}

\begin{passage}{Passage 2}{The following passage is from an opinion piece written in 2025 by a member of a city council for the city's newspaper.}
\pn{1} Somewhere in this city's zoning code, on page 214, is a rule that has shaped our streets more than any architect: every new restaurant must provide one parking space for every 100 square feet of dining room. A coffee shop the size of a living room needs four. A family restaurant needs twenty-five, which is more land for cars than for customers.

\pn{2} Nobody voted for that outcome directly. The rule was copied, in 1961, from a suburban county two states away, and it has been copied forward in every revision since. It was a guess, and it hardened into law.

\pn{3} Its costs are easy to see if you know where to look. A surface parking space here costs about \$9,000 to build; a garage space, closer to \$35,000. Those costs do not disappear. They are folded into rents, into menu prices, and into the price of every apartment that must come with a space whether or not its tenant owns a car. Meanwhile, a count this summer by the planning office found that our downtown lots were, at their busiest hour, 58 percent full.

\pn{4} Opponents of reform worry that without minimums, drivers will circle the block for hours and spill into residential streets. That is a reasonable fear, and in a few neighborhoods it may come true for a while. But the remedy for a parking shortage on one block is to manage that block, with meters or permits, not to require every business in the city to build for a crowd that never arrives.

\pn{5} Ending minimums does not ban parking. A grocer who believes her customers need forty spaces may still build forty. The change simply lets her decide, as she already decides how many registers to install, instead of letting a guess from 1961 decide for her.
\end{passage}

\begin{q}{1}{2.2 Evidence}{Concrete, slightly absurd cases illustrate the rule's effect.}
The examples of the coffee shop and the family restaurant in paragraph 1 primarily serve to
\ch{illustrate the rule's effects with concrete and somewhat absurd consequences}{establish the writer's expertise as a restaurant owner}{show that the rule applies only to large businesses}{prove that most restaurants in the city are losing money}
\end{q}

\begin{q}{1}{2.1 Claims}{The rule began as an unexamined borrowing and persisted by repetition.}
The statement ``It was a guess, and it hardened into law'' (paragraph 2) suggests that the parking rule
\ch{originated without local justification and persisted through repetition}{was approved by a vote of city residents}{was carefully studied before it was adopted}{has been revised to reflect current conditions}
\end{q}

\begin{q}{2}{2.2 Evidence}{Construction costs show how required spaces raise prices passed on to others.}
The construction costs in paragraph 3 are offered mainly to support the claim that
\ch{garages are safer than surface lots}{required parking raises costs that are passed on to residents and customers}{downtown lots are rarely used}{the city should build more parking garages}
\end{q}

\begin{q}{1}{2.2 Evidence}{Lots only 58\% full at peak suggest supply already exceeds demand.}
The planning office's finding that downtown lots were ``at their busiest hour, 58 percent full'' (paragraph 3) functions primarily to
\ch{suggest that existing parking already exceeds demand, undercutting the need for minimums}{show that parking construction is too expensive}{prove that residents prefer walking to driving}{concede that some neighborhoods lack parking}
\end{q}

\begin{q}{3}{2.4 Concession}{He grants the fear is reasonable, then proposes a targeted remedy.}
In paragraph 4, the writer responds to opponents of reform by
\ch{attacking their motives for opposing reform}{dismissing their concerns as unfounded}{granting that their fear is reasonable before proposing a more targeted solution}{agreeing that minimums should remain in residential areas}
\end{q}

\begin{q}{3}{2.4 Qualification}{``A few neighborhoods'' and ``for a while'' limit the concession's scope and duration.}
The phrase ``in a few neighborhoods it may come true for a while'' (paragraph 4) qualifies the writer's concession by
\ch{suggesting the problem will grow worse over time}{admitting that the reform will fail}{limiting both how widely and how long the problem may occur}{blaming residents for the shortage}
\end{q}

\begin{q}{4}{4.3 Comparisons}{Likening parking to registers presents it as an ordinary business choice.}
In paragraph 5, the comparison between deciding on parking spaces and deciding ``how many registers to install'' primarily serves to
\ch{suggest that grocers have too few registers}{criticize grocers for poor planning}{show that parking is more expensive than registers}{present parking as an ordinary business decision owners can make for themselves}
\end{q}

\begin{q}{2}{2.2 Evidence}{Evidence that lots are full would undercut the claim that minimums build for a crowd that never arrives.}
Which finding, if true, would most weaken the writer's argument?
\ch{A garage space costs more than a surface space.}{When the count was repeated in December, downtown lots were over 95 percent full every evening.}{Some grocers prefer to build more parking than the law requires.}{The suburban county that wrote the original rule has since repealed it.}
\end{q}

\section{Passage 3: Personal finance (writing)}

\begin{passage}{Passage 3}{The following is a draft of an argument essay a student is writing for a school publication. It may need revision.}
\pn{1} \sn{1} Most students at our school can solve a quadratic equation. \sn{2} Far fewer could explain what the interest rate on a credit card actually costs them. \sn{3} My cousin, for example, opened her first card at eighteen and spent two years paying off a laptop that ended up costing her nearly twice its price. \sn{4} Our school should require a one-semester personal finance course for graduation.

\pn{2} \sn{5} Financial mistakes made at eighteen can follow a person for years. \sn{6} A low credit score can raise the cost of a car loan, keep a renter from getting an apartment, and even show up in some employers' background checks. \sn{7} Credit cards first became common in the 1950s. \sn{8} Young adults are among the least likely to understand these consequences, often because no one was ever required to teach them.

\pn{3} \sn{9} A required course would fix this problem completely. \sn{10} Researchers have found that graduates in states that added such requirements were more likely to pay their bills on time. \sn{11} The course would not need to be expensive.

\pn{4} \sn{12} Some argue that families, not schools, should teach money management. \sn{13} In an ideal world, they would. \sn{14} But many parents never learned these skills themselves, and a family that is struggling financially may have little time to teach them. \sn{15} A school course gives every student the same starting point.

\pn{5} \sn{16} We already require students to learn the chemistry of baking soda. \sn{17} We can require them to learn the arithmetic of a credit card bill.
\end{passage}

\begin{q}{3}{2.5 Writing: evidence}{Sentence 7's history of credit cards does not support paragraph 2's claim about lasting consequences.}
Which sentence in paragraph 2 should be deleted because it does not support the paragraph's claim?
\chh{Sentence 5}{Sentence 6}{Sentence 7}{Sentence 8}
\end{q}

\begin{q}{2}{2.4 Qualification}{The revision limits the overstatement while keeping the claim.}
The writer wants to revise sentence 9 (reproduced below) to make the claim more defensible.

\emph{A required course would fix this problem completely.}

Which version of sentence 9 best accomplishes this goal?
\ch{A required course would fix this problem completely and permanently.}{A required course would not prevent every mistake, but it would give students the tools to avoid the most common and costly ones.}{Nobody can say whether a required course would help.}{A required course might or might not help.}
\end{q}

\begin{q}{4}{2.5 Writing: evidence}{Naming the study, its date, and the size of the effect makes the finding specific and credible.}
The writer wants to strengthen the evidence in sentence 10. Which addition would best accomplish this goal?
\ch{A sentence noting that some students dislike math}{A quotation from a celebrity about saving money}{A sentence explaining that the researchers were enthusiastic about their work}{The name of the study, when it was published, and how large the improvement was}
\end{q}

\begin{q}{2}{2.5 Writing: evidence}{Existing teachers and an existing elective directly support the low-cost claim.}
The writer wants to add a sentence after sentence 11 to support the claim it makes. Which choice best accomplishes this goal?
\ch{Many students would enjoy learning about money.}{The district's two business teachers already teach an elective on this material and could add sections without new hires.}{Some private schools charge high tuition.}{Textbooks have become more expensive in recent years.}
\end{q}

\begin{q}{3}{2.4 Concession}{Granting part of the opposing view makes the rebuttal that follows more credible.}
The writer is considering deleting sentence 13. Should the writer keep or delete this sentence?
\ch{Delete it, because it abandons the writer's position.}{Delete it, because it repeats sentence 12.}{Keep it, because it concedes the appeal of the opposing view, which makes the rebuttal that follows more credible.}{Keep it, because it provides statistical evidence.}
\end{q}

\begin{q}{1}{2.3 Commentary}{The added sentence explains why an equal starting point matters for the argument.}
The writer wants to add a sentence after sentence 15 that explains why a ``same starting point'' matters to her argument. Which choice best accomplishes this goal?
\ch{Without it, knowing how money works stays a privilege of the families who already have it.}{The course could be taught in the spring semester.}{Starting points are important in races as well.}{Some students already know a great deal about money.}
\end{q}

\begin{q}{4}{2.1 Thesis}{The revised thesis takes the position and previews both main reasons.}
The writer wants to revise sentence 4 so that it previews the essay's main reasons. Which version best accomplishes this goal?
\ch{(as it is now)}{Our school should require a personal finance course, which is something to think about.}{Personal finance is a topic that affects everyone.}{Our school should require a one-semester personal finance course, because early financial mistakes are costly and not every family can teach students to avoid them.}
\end{q}

\begin{q}{2}{4.2 Syntax}{The parallel contrast is memorable; the replacement only restates.}
The writer is considering replacing sentences 16 and 17 with ``In conclusion, a personal finance course is a good idea.'' Should the writer make this change?
\ch{Yes, because a conclusion should restate the thesis directly.}{No, because the original's parallel contrast between a required science fact and an untaught life skill makes the case memorably, while the replacement only restates.}{Yes, because sentences 16 and 17 introduce an unrelated topic.}{No, because the replacement contradicts the thesis.}
\end{q}

\section{Free response}

\begin{frq}{Question 3 \textperiodcentered\ Argument \textperiodcentered\ 6 points \textperiodcentered\ 40 minutes}
\small
\textbf{Directions.} In a commencement address delivered in 2021, the novelist R.\,A.\ Benning (a fictional writer) told graduates: ``Comfort is the enemy of every skill worth having.''

\textbf{Prompt.} Write an essay that argues your position on the extent to which Benning's claim is valid.

In your response you should do the following:
\begin{itemize}[leftmargin=1.5em,itemsep=1pt]
\item Respond to the prompt with a thesis that presents a defensible position.
\item Provide evidence to support your line of reasoning.
\item Explain how the evidence supports your line of reasoning.
\item Use appropriate grammar and punctuation in communicating your argument.
\end{itemize}
\end{frq}

\makekey

\subsection*{Essay scoring guide (6 points)}
\small
The rows follow the published AP English Language rubric (thesis 1, evidence and commentary 4, sophistication 1). The ``earns credit'' column shows what a top response does with this prompt.
\begin{compare}{Row}{Earns credit}{Common miss}
A. Thesis (1) & A defensible position on the claim, e.g.\ ``Benning is right that skill grows only at the edge of discomfort, but she is wrong to call comfort the enemy: the security to fail safely is what lets people keep returning to that edge.'' & ``Comfort can be good or bad''; restating the quotation; a thesis that only defines comfort. \\
B. Evidence and commentary (4) & Two or three specific, developed examples, each tied to a claim in a clear line of reasoning, e.g.\ deliberate practice in music or sport (improvement comes from drilling what one cannot yet do, not repeating what feels easy); an athlete or scientist whose breakthrough followed repeated failure; a personal example (learning a language by speaking it badly in public). Commentary explains \emph{why} each case supports the claim: the mechanism by which discomfort produces growth, or why some comfort (a coach's trust, a safety net) makes risk possible. 4 points: every claim supported and consistently explained. 3: explanation for only some. 2: evidence with little line of reasoning. 1: general or hypothetical evidence, summarized. & Lists of famous names with no detail; extended hypotheticals (``imagine a person who\ldots''); examples narrated without explaining how they prove the point. \\
C. Sophistication (1) & Explores the tension in the prompt: who can afford discomfort (financial security, privilege), the difference between productive struggle and harm, or how comfort, as rest and safety, enables the next risk; or situates the claim in a broader context (training science, education policy); or sustains a vivid, controlled style. & A last-sentence twist that the body never developed; sweeping generalizations about ``society today.'' \\
\end{compare}

\end{document}
